\chapter{Adult Jaundice}

\section*{Definition}
Adult jaundice is yellow discoloration of the skin or whites of the eyes caused by increased bilirubin in the blood. It is a sign rather than a diagnosis. Causes range from hemolysis and liver inflammation to blockage of bile flow, and new jaundice requires medical evaluation.

\section*{What Happens in Your Body}
Red blood cells produce bilirubin as hemoglobin is broken down. The liver processes bilirubin and excretes it into bile. Jaundice develops when bilirubin production increases, liver processing fails, or bile cannot drain. Conjugated bilirubin may darken urine, while impaired bile delivery can make stools pale and cause itching.

\section*{Causes}
\begin{itemize}
  \item Viral, autoimmune, alcoholic, metabolic, or drug-induced hepatitis
  \item Gallstones, bile-duct obstruction, pancreatitis, or biliary cancer
  \item Hemolysis, ineffective red-cell production, or inherited bilirubin disorders
  \item Cirrhosis, sepsis, or acute liver failure
  \item Liver, pancreatic, gallbladder, or bile-duct tumors
\end{itemize}

\section*{When to See a Doctor}
Seek urgent care for jaundice with confusion, severe sleepiness, bleeding, fever, chills, severe abdominal pain, repeated vomiting, rapidly worsening symptoms, or very dark urine and pale stools. Jaundice during pregnancy, after a new medication, or with acetaminophen overdose requires prompt assessment.

\section*{Related Symptoms}
\begin{itemize}
  \item Dark urine, pale or clay-colored stool, itching, nausea, and loss of appetite
  \item Right-upper abdominal pain, fever, weight loss, or abdominal swelling
  \item Easy bruising, bleeding, confusion, or drowsiness in advanced liver dysfunction
\end{itemize}

\section*{Self-Care / Home Management}
Avoid alcohol and non-prescribed supplements or medications until evaluated. Do not exceed labeled acetaminophen doses. Maintain hydration if able to drink, but do not delay care for new jaundice or attempt a detoxification regimen.

\section*{How Your Doctor Will Evaluate You}
\begin{itemize}
  \item History addresses alcohol, medication and supplement exposure, travel, infection risk, transfusion, family history, pregnancy, and stool or urine changes.
  \item Examination assesses liver and spleen enlargement, abdominal tenderness, ascites, bruising, fever, and neurologic status.
  \item Tests include fractionated bilirubin, liver enzymes, alkaline phosphatase, albumin, INR, blood count, hemolysis studies, viral testing, and ultrasound; CT or MRCP may follow.
\end{itemize}

\section*{Treatment Options}
Treatment targets the cause: stopping a toxic drug, treating infection or autoimmune disease, relieving obstruction, managing hemolysis, or providing specialist care for acute liver failure. Biliary infection or obstruction may require urgent endoscopy or surgery.

\section*{References}
\begin{thebibliography}{99}
\bibitem{adult_jaundice:ref1} Kwo PY, Cohen SM, Lim JK. ACG Clinical Guideline: Evaluation of Abnormal Liver Chemistries. \textit{Am J Gastroenterol}. 2017;112:18--35.
\bibitem{adult_jaundice:ref2} European Association for the Study of the Liver. EASL Clinical Practice Guidelines on acute liver failure. \textit{J Hepatol}. 2017;66:1047--1081.
\bibitem{adult_jaundice:ref3} American College of Radiology. ACR Appropriateness Criteria: Jaundice. Revised 2023.
\end{thebibliography}
